\subsection{自由李代数的 Hall 基}

我们沿用上一 no. 的记号。

定理 1. 设 H 是相对于 X 的一个 Hall 集， \(\Psi\) 是 M(X) 到自由李代数 L(X) 中的典范映射。 \(\Psi\) 在 H 上的限制是模 L(X) 的一个基。

对 H 的每个元素 \(w\) ，我们记 \(\bar{w} = \Psi(w)\) 。

\begin{enumerate}
\def\labelenumi{(\Alph{enumi})}
\tightlist
\item
  X 为有限集的情形。
\end{enumerate}

若 X 是空集，则 \(M(X)\) 也是空集，因而 H 也是，并且 \(L(X)\) 为零。若 X 只有一个元素 x，则 \(H \cap M^{n}(X)\) 是空集（对 \(n \geqslant 2\) ）(命题 12 (a))。因此 \(H = \{x\}\) ；我们还知道 (no. 2， 注)，模 \(L(X)\) 是自由的并且以 \(\{\bar{x}\}\) 为基。因此当 X 至多有一个元素时定理成立。

以下设 X 至少有两个元素；选取具有命题 13 所述性质的序列 \((w_{p})\) 与 \((P_{p})\) 。对每个整数 \(p \geqslant 0\) ，用 \(L_{p}\) 表示 \(\mathrm{L}(\mathbf{X})\) 中由元素 \(\bar{w}_{i}\) （满足 \(0 \leqslant i < p\) ）生成的子模，并用 \(g_{p}\) 表示 \(\mathrm{L}(\mathbf{X})\) 中由族 \((\bar{u})_{u \in \mathbf{P}_{p}}\) 生成的李子代数。

引理 2. 对每个整数 \(p \geqslant 0\) ，模 \(\mathrm{L}_p\) 以族 \((\bar{w}_i)_{0 \leqslant i < p}\) 为基，李代数 \(\mathfrak{g}_p\) 以 \((\bar{u})_{u \in \mathbf{P}_p}\) 为基本族，并且模 \(\mathrm{L}(\mathbf{X})\) 是 \(\mathrm{L}_p\) 与 \(\mathfrak{g}_p\) 的直和。

\(L_{0} = \{0\}\) 且 \(g_{0} = L(X)\) ，因而引理对 \(p = 0\) 成立。我们对 \(p\) 作归纳论证。设引理对某个整数 \(p \geqslant 0\) 成立。令 \(u_{i,w} = (\text{ad } \bar{w}_{p})^{i}. \bar{w} = \Psi(w_{p}^{i} w)\) （对每个 \(i \geqslant 0, w \in P_{p}, w \neq w_{p}\) ）。由 no. 9 命题 10 的推论，自由李代数 \(g_{p}\) 是模 \(T_{p}\) （以 \(\{\bar{w}_{p}\}\) 为基）与一个李子代数 \(b_{p}\) 的直和，后者以

\[
\mathcal {F} = (u _ {i, w}) _ {i \geqslant 0, w \in \mathrm{P} _ {p}, w \neq w _ {p}}
\]

为基本族的李子代数的直和。由命题 13 (d)，族 \((\bar{u})_{u\in \mathbf{P}_{p + 1}}\) 等于 \(\mathcal{F}\) ，因而是 \(\mathfrak{h}_p = \mathfrak{g}_{p + 1}\) 的一个基本族。因此 \(\mathrm{L}(\mathrm{X}) = \mathrm{L}_p\oplus \mathrm{T}_p\oplus \mathfrak{g}_{p + 1}\) ，并且由于 \(\mathrm{L}_{p + 1} = \mathrm{L}_p + \mathrm{T}_p\) ，有 \(\mathrm{L}(\mathrm{X}) = \mathrm{L}_{p + 1}\oplus \mathfrak{g}_{p + 1}\) ，而 \((\bar{w}_0,\bar{w}_1,\dots ,\bar{w}_{p - 1},\bar{w}_p)\) 是模 \(\mathbf{L}_{p + 1}\) 的一个基。

设 \(n\) 是一个正整数。由命题 13 (c)，存在一个整数 \(p(n)\) ，使得 \(P_p\) 只含有长度为 \(>n\) 的元素，其中 \(p \geqslant p(n)\) 。对 \(p \geqslant p(n)\) ，李子代数 \(g_p\) （\(L(X)\) 中的）由 \(>n\) 次的元素生成，因此 \(L^n(X) \cap g_p = \{0\}\) 。另一方面，元素 \(\bar{w}_i\) （\(L(X)\) 中的）都是齐次的，并且族 \((\bar{w}_i)_{0 \leqslant i < p}\) 是 \(g_p\) 的一个补模的基。由此立即得出，由元素 \(\bar{w}_i\) （次数为 \(n\) ）组成的族是模 \(L^n(X)\) 的一个基，并且序列 \((\bar{w}_i)_{i \geqslant 0}\) 是模 \(L(X)\) 的一个基。

\begin{enumerate}
\def\labelenumi{(\Alph{enumi})}
\setcounter{enumi}{1}
\tightlist
\item
  一般情形。
\end{enumerate}

若 \(S\) 是 \(X\) 的一个子集，回顾 \(M(S)\) 与 \(M(X)\) 的这样一个子 magma 等同：它由 \(S\) 生成；而 \(L(S)\) 与 \(L(X)\) 的这样一个李子代数等同：它由 \(S\) 生成；我们已经看到，若 \(w \in M(S)\) 的长度为 \(\geqslant 2\) ，则 \(\alpha(w) \in M(S)\) 且 \(\beta(w) \in M(S)\) 。由此立即得出 \(H \cap M(S)\) 是相对于 \(S\) 的 Hall 集。

对每个有限子集 \(\Phi\) （\(H\) 的），存在一个有限子集 \(S\) （\(X\) 的）使得 \(\Phi \subset \mathrm{M}(\mathrm{S})\) 。于是情形 (A) 表明，元素 \(\bar{w}\) （带有 \(w \in \Phi\) ）在 \(\mathrm{L}(\mathrm{S})\) 中线性无关，因而在 \(\mathrm{L}(\mathrm{X})\) 中也线性无关。因此族 \((\bar{w})_{w \in \mathbf{H}}\) 是自由的。

对每个元素 \(a\) （\(\mathrm{L}(\mathrm{X})\) 中的），存在一个有限子集 \(\mathrm{S}\) （\(\mathrm{X}\) 的）使得 \(a \in \mathrm{L}(\mathrm{S})\) 。由情形 (A)，子集 \(\Psi(\mathrm{H} \cap \mathrm{M}(\mathrm{S}))\) （\(\Psi(\mathrm{H})\) 中的）生成模 \(\mathrm{L}(\mathrm{S})\) ，因而 \(a\) 是 \(\Psi(\mathrm{H})\) 的元素的线性组合。因此 \(\Psi(\mathrm{H})\) 生成模 \(\mathrm{L}(\mathrm{X})\) ，这就完成了证明。

推论. 模 \(\mathbf{L}(\mathbf{X})\) 是自由的，并且子模 \(\mathbf{L}^{\alpha}(\mathbf{X})\) （其中 \(\alpha \in \mathbf{N}^{(\mathbf{X})}\) ）以及 \(\mathbf{L}^n (\mathbf{X})\) （其中 \(n\in \mathbf{N}\) ）也都是自由的。模 \(\mathbf{L}^{\alpha}(\mathbf{X})\) 都是有限秩的，并且模 \(\mathbf{L}^n (\mathbf{X})\) 也都是有限秩的（若 \(\mathbf{X}\) 有限）。

存在相对于 X 的 Hall 集 H (命题 11)。对一切 \(w \in \mathbf{H}\) ，元素 \(\Psi(w)\) （\(\mathbf{L}(\mathbf{X})\) 中的）属于诸模 \(\mathrm{L}^{\alpha}(\mathbf{X})\) 之一（其中 \(\alpha \in \mathbf{N}^{(\mathbf{X})}\) ），而模 \(\mathrm{L}(\mathbf{X})\) 是诸子模 \(\mathrm{L}^{\alpha}(\mathbf{X})\) 的直和。此外，对一切 \(\alpha \in \mathbf{N}^{(\mathbf{X})}\) ， \(\mathrm{M}(\mathbf{X})\) 中其典范像在 \(\mathbf{N}^{(\mathbf{X})}\) 中等于 \(\alpha\) 的元素组成的集合是有限的；这表明诸模 \(\mathrm{L}^{\alpha}(\mathbf{X})\) 中的每一个都是自由且有限秩的，并且 \(\mathrm{L}(\mathbf{X})\) 是自由的。而 \(\mathrm{L}^{n}(\mathbf{X}) = \bigoplus_{|\alpha| = n} \mathrm{L}^{\alpha}(\mathbf{X})\) ，因而 \(\mathrm{L}^{n}(\mathbf{X})\) 是自由的；当 X 有限时，由 \(\alpha \in \mathbf{N}^{(\mathbf{X})}\) （满足 \(|\alpha| = n\) ）组成的集合是有限的，因而 \(\mathrm{L}^{n}(\mathbf{X})\) 是有限秩的。

定义 3. 自由李代数 \(\mathbf{L}(\mathbf{X})\) 的 Hall 基是指 \(\mathbf{L}(\mathbf{X})\) 的任一这样的基：它是一个相对于 \(\mathbf{X}\) 的 Hall 集的典范像。

注. 假设 X 由两个不同元素 x 与 y 组成，并设 \(L^{(\cdot,1)}\) 是 \(L(X)\) 的由诸 \(L^{\alpha}(X)\) （其中 \(\alpha\in N^{X}\) 且 \(\alpha(y)=1\) ）之和组成的子模。由 no. 10 的定理 1 与命题 12 立即得出， \((ad x)^{n}\cdot y\) （其中 n 是一个整数 \(\geqslant0\) ）的元素构成子模 \(L^{(\cdot,1)}\) 的一个基。由此得出，映射 ad x 在 \(L^{(\cdot,1)}\) 上的限制是单射。

